\documentclass[a4paper, 12pt]{book}

\usepackage{graphicx}
\usepackage{wrapfig}
\usepackage[font=small,labelfont=bf]{caption}
\usepackage{index}
\usepackage{hyperref}
\usepackage{xcolor}
\usepackage{listings}
\usepackage{xepersian}

\makeatletter
\renewcommand\@endpart{\vfil
              \if@twoside
                \null
                \thispagestyle{empty}%
                \newpage
              \fi
              \if@tempswa
                \twocolumn
              \fi}
\makeatother

\renewcommand{\baselinestretch}{1.6}

\newcommand{\lrbold}[1]{\lr{\textbf{#1}}}
\newcommand{\lrboldit}[1]{\lr{\textbf{\textit{#1}}}}
\newcommand{\lrit}[1]{\lr{\textit{#1}}}


\settextfont{B Yekan}
\setlatintextfont[Scale=1.2]{Arial}
\DefaultMathsDigits

% -- define code style 

\definecolor{customblue}{RGB}{59, 59, 59}
\definecolor{keywordcolor}{RGB}{255, 196, 0}
\definecolor{codepurple}{rgb}{0.58,0,0.82}
\definecolor{stringcolor}{RGB}{255, 196, 0}
\definecolor{light-gray}{gray}{0.95}
\lstdefinestyle{C++Style}{%
  backgroundcolor=\color{customblue},
  breaklines=true,
  basicstyle=\footnotesize\ttfamily,
  keywordstyle=\color{keywordcolor},
  commentstyle=\color{stringcolor}\textit,
  stringstyle=\color{codepurple},
  numbers=left,
  numberstyle={\tiny\lr},
  showspaces = false,
  showstringspaces = false,
  tabsize = 2,
  frame=single,
  xleftmargin=5pt,
  xrightmargin=3pt,
  language =  C++,
  aboveskip = 20pt,
  rulecolor=\color{stringcolor},
  captiondirection=RTL,
}


\lstnewenvironment{C++Code}[1][]
{%
    \lstset{style=C++Style, #1}%
}{}


% end code style defenition

\title{ مقدمه ای بر گرافیک کامپیوتر}
\author{احمد منصوری و \lr{\textsf{\textbf{peter shirley}}}}
\date{\lr{December 21, 2009}}

\begin{document}
\maketitle
\let\cleardoublepage\clearpage

\huge\textbf{چکیده}
\normalsize
\begin{flushright}
  همزمان با پیدایش کامپیوتر ها، تلاش ها برای بهره بردن از توان آنها برای ابزار های \lr{Visualize} و دیگر ابزار ها برای
  استفاده از این قابلیت در زمینه های نظامی، فیلم و انیمیشن، شبیه سازی و بازی سازی شروع شد، این ابزار ها یا به صورت
  اختصاصی برای استفاده در صنایع خاص طراحی می شوند یا به صورت عمومی تر برای کاربرد های وسیع تری طراحی و توسعه می یابند.
   طراحی و پیاده سازی موتور های گرافیکی به صورت کلی دارای پایه ها و دانشی یکسان از نحوه کار پردازنده ها و ریاضیات است.

\end{flushright}


\begingroup
  \hypersetup{hidelinks}
  \tableofcontents
\endgroup

\makeatletter
\renewcommand\@endpart{\vfil
              \if@twoside
                \null
                \thispagestyle{empty}%
                \newpage
              \fi
              \if@tempswa
                \twocolumn
              \fi}
\makeatother

\part{گرافیک}

\huge
    مقدمه \\

\vspace*{2cm}
\noindent
\normalsize
     در بخش اول به معرفی و توضیح قسمت های مربوط به \textbf{تصویر} در پروژه می پردازم،
     این قسمت از موتور مسئولیت دریافت مدل های سه بعدی و اطلاعات مربوطه(\lr{\textbf{Textures}, \textbf{Normals}, \textbf{geometry}, ...})
     و نمایش آن ها در صفحه را برعهده دارد، در این قسمت ما با استفاده از ریاضیات مدل ها را در فضای سه بعدی شبیه سازی می کنیم.\par
     تمامی ابزار های استفاده شده در این  قسمت، در طول فصول و متناسب با بخشی که از آن ها استفاده شده معرفی می شوند.\par
     هر فصل در این بخش مستقیما مربوط به یکی از قسمت های موتور در بخش گرافیک است، ابتدای هر فصل فایل های مربوطه به آن فصل ذکر خواهند شد.

\chapter{\lr{Renderer}}
\newpage

\addcontentsline{toc}{section}{پنجره ها و کانتکست ها}
\section*{\huge{پنجره ها و کانتکست ها}}
\vspace*{0.6cm}

\subsection*{\lr{openGL}}
\noindent
\normalsize
    \lrboldit{OpenGL} یک \lrboldit{api} چند زبانه و کراس پلتفرم است که برای به تصویر کشیدن تصاویر دو بعدی و سه بعدی با استفاده از بردار ها
    اسفاده می شود، معمولا از \lrboldit{OpenGL} برای برقراری ارتباط با واحد پردازش گرافیکی \lrboldit{GPU} و بهره بردن از سرعت سخت افزار مخصوص برای رندر استفاده می شود.\par
    
    همچنین \lrit{api} های دیگری نیز برای استفاده از قدرت سخت افزاری و پردازنده گرافیکی وجود دارند \lrit{Vulkan} نیز مانند \lrit{openGL} چند زبانی و چند سکویی است،
    \lrit{DirectX} به صورت انحصاری توسط مایکروسافت توسعه می یابد و در سیستم عامل ویندوز استفاده می شود ، شرکت \lrbold{Apple} از \lrit{api} اختصاصی خود به نام \lrit{Metal}
    به صورت انحصاری پشتیبانی می کند، دلیل انتخاب\lrit{openGL} در این پروژه چندسکویی بودن و ساختار ساده تر برای پروژه های آموزشی در زمینه \lrit{real-time rendering} می باشد.\par
    
    نباید \lrit{OpenGL} را با یک کتابخانه\lrit{(library)} اشتباه گرفت، \lrbold{OpenGL} به صورت یک \lr{interface} و یک قراداد انتزاعی در ورژن های مختلف ارائه می شود که فروشندگان و سازندگان\lrit{(vendor)} مختلف باید پیاده سازی ای منطبق با این قرارداد را انجام دهند.
    پس از نصب درایور مربوط به پردازنده گرافیکی، برنامه نویس قابلیت دسترسی به تابع های مختلف که توسط \lrit{vendor} پیاده سازی شده را خواهد داشت.
    برای استفاده از \lrit{\textsf{OpenGl}} نیاز به ابزار های دیگری نیز داریم، ابتدا نیاز داریم که یک \lrit{window} و یک \lrit{context} تعریف کنیم،
    برای این کار از \lrbold{GLFW} استفاده می کنیم.


\newpage
\subsection*{\lr{GLFW}}
\noindent
\normalsize
    یک کتابخانه برای ساختن \lrit{window} و \lrit{context} ها برای \lrit{openGL, openGL ES, Vulkan} است، این کتابخانه به زبان \lrbold{C} نوشته شده و \lrit{binding} های مختلف آن به زبان های مختلف موجود است، این کتابخانه همچنین توانایی کنترل کردن ورودی های مختلف مثل \lrit{keyboard, mouse, joystick} را داراست، ما برای استفاده از \lrit{openGL} نیاز به این کتابخانه یا مشابه آن داریم زیرا \lr{openGL} هیچگونه قابلیت پیشفرضی برای مدیریت \lr{window} یا \lr{context} ها یا مدیریت \lr{input} ندارد. همچنین \lr{GLFW} یک کتابخانه چندسکویی است و می توانیم آن را در سیستم عامل های مختلف استفاده کنیم، پروژه من نیز چندسکویی است، پس می توانیم از این کتابخانه سبک و چندسکویی استفاده کنیم.
    \lr{windows}، پنجره ای است که \lr{GLFW} به وسیله امکانات فراهم شده در سطح سیستم عامل برای ما فراهم می کند، همچنین یک \lr{context} را می توانیم به عنوان یک شئ در نظر بگیریم که تمامی اطلاعات \lr{openGL} را به همراه دارد، اطلاعاتی مانند \lr{state} و  \lr{framebuffers}ها .
    برای کنترل کردن ورودی ها، \lr{glfw} از دو روش استفاده می کند، برای ورودی\lrit{mouse} از \lr{callback function} ها استفاده می کند، اما برای ورودی \lr{keyboard} می توانیم از تابع های کتابخانه استفاده کنیم و به صورت مستقیم ورودی را دریافت کنیم.\par
    حالا که به \lr{window} و \lr{context} دسترسی داریم، باید دسترسی به تابع های \lr{opengl} فراهم کنیم، برای این کار از \lrbold{GLAD} استفاده می کنیم.

\begin{figure}[ht]
    \centering
    \href{https://github.com/glfw}{
        \includegraphics[width=3cm]{GLFW.png}
    }
    \caption{\lr{\textit{glfw logo}}}
    \label{fig:my_label}
\end{figure}

\newpage
\subsection*{\lr{GLAD}}
\noindent
\normalsize
    کتابخانه ای برای \lrit{load} کردن \lr{pointer} ها به توابع \lr{opengl} در هنگام \lr{runtime}.
    این کتابخانه یکی از کتابخانه های\lr{OpenGL Loading Library} است، برای کار با \lr{opengl} ما حتما باید یکی از این کتابخانه هارا مورداستفاده قرار دهیم تا بتوانیم به توابع \lr{opengl} دسترسی داشته باشیم، این کتابخانه ها هم ویژگی های \lr{Core} که توسط \lr{opengl} مشخص شده را \lr{load} می کنند و هم ویژگی های \lr{extension} که توسط \lr{Vendor} ها به پیاده سازی آن ها از \lr{opengl} اضافه شده، علاوه بر این دیگر نیازی به اضافه کردن فایل های مربوط به \lr{opengl} نیست و این فایل ها به صورت خودکار همه موارد را تنطیم می کنند.
    \lr{Glad} یک \lr{generator} است که براساس پارامتر هایی که کاربر انتخاب می کند یک فایل حاوی تمامی تعریف های مربوط به \lr{constant} و تابع ها و ... به ما ارائه می کند، بعد از دانلود این فایل و اضافه کردن به آن به پروژه از طریق کد زیر می توانیم تمامی \lr{opengl function poiter} هارا در \lr{runtime} بارگزاری کنیم.
    
    \begin{LTR}
        \begin{lstlisting}[style=C++Style,caption=\rl{load opengl function pointer}]
          int main()
          {
            return 0;
          }
        \end{lstlisting}
    \end{LTR}


\part{فیزیک}
\begin{center}
    \chapter{\lr{Physics}}
\end{center}
\section{\lr{Particles}}
\section{\lr{Rigid Bodies}}

\end{document} 